\begin{abstract_en}

	Appropriate mathematical models can accurately describe the dynamic characteristics of specific epidemics.
	This paper aims to explore new approaches in epidemiological research by integrating epidemiological investigations, differential equation models, and neural networks to more accurately predict and monitor the transmission trends of infectious diseases.

	In this paper, we establish an SEIRV model with age structure and vaccination and LSTM model.
	We also discuss the application of epidemiological investigation methods, differential equation models, and neural networks in addressing issues related to epidemiological dynamics.
	Through practical case studies, we validate the effectiveness of this integrated approach.

	The research findings indicate that the combination of differential equation models and neural networks holds immense potential in infectious disease modeling.
	This approach can capture transmission dynamics more accurately and predict future trends, providing powerful tools for combating infectious disease outbreaks.
\end{abstract_en}

\begin{keywords_en}
	Infectious Disease Modeling;
	Epidemiology;
	Differential Equations;
	Neural Networks;
	COVID-19
\end{keywords_en}